$D = \C$: $f^{2} = -1$ means $f$ is annihilated by $x^{2} + 1$, which has distinct roots $\pm\mathrm{i}$, so $f$ is diagonalisable with eigenvalues in $\{\mathrm{i}, -\mathrm{i}\}$; it is conjugate to $\operatorname{diag}(\mathrm{i} 1_k, -\mathrm{i} 1_{m-k})$ with $k$ the dimension of the $\mathrm{i}$-eigenspace, which is a conjugation invariant.

$D = \R$: $f$ makes $\R^{m}$ a complex vector space with $\mathrm{i}$ acting as $f$, so $m = 2 \dim_\C$ is even; choosing a complex basis $b_1, \dots, b_{m/2}$, the real basis $b_1, f b_1, \dots$ puts $f$ in the standard block form $\operatorname{diag}(J, \dots, J)$ with $J = \begin{psmallmatrix} 0 & -1 \\ 1 & 0 \end{psmallmatrix}$. So all such $f$ are conjugate.

$D = \HH$: let $f$ act $\HH$-linearly on the right $\HH$-module $\HH^{m}$, and let $R_q$ be right multiplication by $q$. Then $f$ commutes with $R_{\mathrm{i}}$ and $(f R_{\mathrm{i}})^{2} = f^{2} R_{\mathrm{i}}^{2} = 1$, so $\HH^{m} = E_+ \oplus E_-$ with $E_\pm$ the $\pm 1$-eigenspaces of $f R_{\mathrm{i}}$. $R_{\mathrm{j}}$ commutes with $f$ and anticommutes with $R_{\mathrm{i}}$, so it exchanges $E_+$ and $E_-$, and $E_+$ is invariant under $R_{\mathrm{i}}$, a complex space of dimension $m$. A complex basis $b_1, \dots, b_m$ of $E_+$ is a quaternionic basis of $\HH^{m}$ (it spans $E_+ \oplus E_+ \mathrm{j} = E_+ \oplus E_-$), and on $E_+$, $f(b_k) \mathrm{i} = b_k$, so $f(b_k) = -b_k \mathrm{i}$. In this basis $f(\sum_k b_k q_k) = \sum_k b_k(-\mathrm{i}\, q_k)$: $f$ is left multiplication by $-\mathrm{i}$ on coordinates, $f = -\mathrm{i}\, 1_m$ as a matrix, the same for every $f$; so all $f$ are conjugate.

Sums: conjugation in $A \oplus A$ acts componentwise. The consequences follow from the table of \claimref{std-algebra.periodicity}: the real types $\Mat_m(\R)$ and $\Mat_m(\R)^{2}$ have $m = 1$ only for $\Cl(0,0) = \R$ and $\Cl(1,0) = \R^{2}$; $\Cl(1,3) \cong \Mat_2(\HH)$ and $I^{2} = -1$; in $\Cl(3,0)$ the centre is spanned by $1$ and $\omega$ (\claimref{std-algebra.complexification}), so $f \mapsto \langle f \omega^{-1}\rangle_0$ is a conjugation invariant, proportional to the complex trace's imaginary part, taking the values $1, -1, 0$ for $k = 2, 0, 1$.
